% Created 2024-01-11 jue 01:46
% Intended LaTeX compiler: pdflatex
\documentclass[11pt]{article}
\usepackage[utf8]{inputenc}
\usepackage[T1]{fontenc}
\usepackage{graphicx}
\usepackage{longtable}
\usepackage{wrapfig}
\usepackage{rotating}
\usepackage[normalem]{ulem}
\usepackage{amsmath}
\usepackage{amssymb}
\usepackage{capt-of}
\usepackage{hyperref}
\usepackage{minted}

\usepackage{parskip}
\usepackage[utf8]{inputenc} %% For unicode chars
\usepackage{placeins}
\usepackage[margin=2.50cm]{geometry}
\usepackage[T1]{fontenc}
\usepackage{mathpazo}
\linespread{1.05}
\usepackage[scaled]{helvet}
\usepackage{courier}
\hypersetup{colorlinks=true,linkcolor=blue}
\RequirePackage{fancyvrb}
\DefineVerbatimEnvironment{verbatim}{Verbatim}{fontsize=\small,formatcom = {\color[rgb]{0.5,0,0}}}
\usepackage{mdframed}
\BeforeBeginEnvironment{minted}{\begin{mdframed}}
\AfterEndEnvironment{minted}{\end{mdframed}}
\setcounter{secnumdepth}{3}
\date{}
\title{}
\hypersetup{
 pdfauthor={},
 pdftitle={},
 pdfkeywords={},
 pdfsubject={},
 pdfcreator={Emacs 28.1 (Org mode 9.6.12)}, 
 pdflang={Spanish}}
\begin{document}

\setcounter{tocdepth}{3}
\tableofcontents

\setlength\parindent{10pt}

\section{PHP Arrays}
\label{sec:orgabfc312}
Arrays are one of the fundamental data structures in PHP, used to store
collections of values. They are versatile and can be used to store
various data types, including numbers, strings, objects, and even other
arrays.

\subsection{Indexed Arrays}
\label{sec:orgd2a9c9b}
\textbf{Creating Arrays}

There are two primary ways to create arrays in PHP:

\begin{enumerate}
\item \textbf{Using the square bracket syntax:} This method involves enclosing the
elements within square brackets, separated by commas.
\end{enumerate}

\begin{minted}[]{php}
<?=
 $numbers = [1, 2, 3, 4, 5];
 $colors = ["red", "green", "blue"];
 $mixed = [1, "hello", true, [1, 2, 3]];
\end{minted}

\begin{enumerate}
\item \textbf{Using the array() function:} This method explicitly creates an array
and assigns it to a variable.
\end{enumerate}

\begin{minted}[]{php}
<?=
 $names = array("Alice", "Bob", "Charlie");
\end{minted}

\textbf{Accessing Array Elements}

To access an element in an array, you can use its index, which is the
position of the element within the array. Indexes start from 0, and you
can access elements using the square bracket syntax.

\begin{minted}[]{php}
<?=
 $firstNumber = $numbers[0]; // Output: 1
 $secondColor = $colors[1]; // Output: green
 $mixedElement = $mixed[3][2]; // Output: 3
\end{minted}

\textbf{Modifying Array Elements}

You can modify the value of an existing element in an array by assigning
a new value to its index.

\begin{minted}[]{php}
<?=
 $names[0] = "John"; // Change the first name to John
 $colors[2] = "purple"; // Change the third color to purple
\end{minted}

\textbf{Adding Elements to Arrays}

To add new elements to an array, you can use the following methods:

\begin{enumerate}
\item \textbf{Using the push() method:} This method appends the new element to the
end of the array.
\begin{minted}[]{php}
<?=
 $names[] = "David"; // Add David to the end of the names array
 $numbers[] = 6; // Add 6 to the end of the numbers array
\end{minted}

\item \textbf{Using the array\_push() function:} This method works similarly to the
push() method but takes an array as the first argument and the
elements to be added as the second argument.
\begin{minted}[]{php}
<?=
 array_push($colors, "yellow", "orange");
\end{minted}

\item \textbf{Using the array\_merge() function:} This method merges two or more
arrays into one.
\begin{minted}[]{php}
<?=
 $newColors = array_merge($colors, ["pink", "cyan"]);
\end{minted}
\end{enumerate}

\textbf{Removing Elements from Arrays}

To remove elements from an array, you can use the following methods:

\begin{enumerate}
\item \textbf{Using the unset() function:} This method removes the element at the
specified index from the array.
\begin{minted}[]{php}
<?=
 unset($numbers[4]); // Remove the fourth element from the numbers array
\end{minted}

\item \textbf{Using the pop() method:} This method removes the last element from
the array and returns its value.
\begin{minted}[]{php}
<?=
 $lastColor = array_pop($colors); // Remove the last color and store it in lastColor variable
\end{minted}

\item \textbf{Using the array\_splice() function:} This method removes elements
from the array and can also replace them with new elements.
\begin{minted}[]{php}
<?=
 array_splice($colors, 1, 1, ["teal", "magenta"]); // Remove the second element and replace it with "teal" and "magenta"
\end{minted}
\end{enumerate}

\textbf{Other Array Operations}
\begin{enumerate}
\item \textbf{count()}.
The \texttt{count()} function is a built-in function in PHP that returns the
number of elements in an array. It is a useful function for
determining the size of an array. Here is an example of how to use the
\texttt{count()} function:
\begin{minted}[]{php}
<?=
  $numbers = [1, 2, 3, 4, 5];
  $numberOfElements = count($numbers);
  echo "Number of elements: $numberOfElements";
\end{minted}

\item \textbf{array\_shift()}
The \texttt{array\_shift()} function removes the first element from an array and
returns it. It is a convenient way to remove and access the first
element of an array. Here is an example of how to use the
\texttt{array\_shift()} function:
\begin{minted}[]{php}
<?=
 $numbers = [1, 2, 3, 4, 5];
 $firstNumber = array_shift($numbers);
 echo "First number: $firstNumber";
\end{minted}

\item \textbf{array\_unshift()}
The \texttt{array\_unshift()} function adds a new element to the beginning of an
array. It is a simple way to insert an element at the beginning of an
array. Here is an example of how to use the \texttt{array\_unshift()} function:
\begin{minted}[]{php}
<?=
 $numbers = [1, 2, 3, 4, 5];
 array_unshift($numbers, 0);
 echo "New array: $numbers";
\end{minted}

\item \textbf{array\_slice()}
The \texttt{array\_slice()} function extracts a specified portion of an array
and returns it as a new array. It is a powerful tool for retrieving
subsets of an array. Here is an example of how to use the
\texttt{array\_slice()} function:
\begin{minted}[]{php}
<?=
 $numbers = [1, 2, 3, 4, 5];
 $slicedNumbers = array_slice($numbers, 2, 2);
 echo "Sliced array: $slicedNumbers";
\end{minted}
\end{enumerate}

These functions are essential for working with arrays in PHP. They allow
you to access, modify, and manipulate arrays effectively.

Arrays are versatile and essential data structures in PHP, allowing you
to store, organize, and manipulate collections of data efficiently.

\subsection{Multidimensional arrays}
\label{sec:orga4bb55c}
Multidimensional arrays in PHP are arrays that contain other arrays as
elements. This allows you to store and organize data in a hierarchical
structure, making it well-suited for representing data with nested
relationships.

\textbf{Creating Multidimensional Arrays}

Multidimensional arrays can be created using the square bracket syntax,
nesting one or more arrays within another. For example, the following
code creates a 2D array:

\begin{minted}[]{php}
<?=
 $students = [
     ["Alice", "12345", "English", "History", "Math"],
     ["Bob", "67890", "Science", "Art", "Computer Science"]
 ];
\end{minted}

In this example, the \texttt{students} array contains two elements, each
representing a student's information. Each student's information is
stored in another array within the \texttt{students} array.

\textbf{Accessing Elements in Multidimensional Arrays}

To access an element in a multidimensional array, you can use nested
indexing, where each index corresponds to the position of the element in
the corresponding subarray. For example, the following code accesses the
first student's name from the \texttt{students} array:

\begin{minted}[]{php}
<?=
 $firstName = $students[0][0];
 echo "First name: $firstName";
\end{minted}

\textbf{Modifying Elements in Multidimensional Arrays}

You can modify an element in a multidimensional array by assigning a new
value to its index. For example, the following code changes the second
student's math grade to 100:

\begin{minted}[]{php}
<?=
 $students[1][4] = 100;
 echo "Second student's math grade: " . $students[1][4];
\end{minted}

\textbf{Adding Elements to Multidimensional Arrays}

To add new elements to a multidimensional array, you can use the same
methods as with one-dimensional arrays. For instance, the following code
adds a new student to the \texttt{students} array:

\begin{minted}[]{php}
<?=
 $students[] = ["Charlie", "98765", "Spanish", "Economics", "Music"];
\end{minted}

\textbf{Removing Elements from Multidimensional Arrays}

To remove elements from a multidimensional array, you can use the
\texttt{unset()} function, followed by the nested indices of the element to be
removed. For example, the following code removes the first student's
information from the \texttt{students} array:

\begin{minted}[]{php}
<?=
 unset($students[0]);
\end{minted}

Multidimensional arrays offer a powerful way to organize and store
complex data structures in PHP. They are particularly useful for
representing nested relationships, such as employee hierarchies, product
categories, or network topologies.

\subsection{Associative arrays}
\label{sec:org42ec974}
Associative arrays in PHP are arrays whose elements are associated with
keys instead of numeric indices. This allows you to store and retrieve
data based on meaningful labels rather than numerical positions.

\textbf{Creating Associative Arrays}

Associative arrays can be created using the same square bracket syntax
as regular arrays, but instead of using numeric indices, you use
key-value pairs enclosed in curly braces. For example, the following
code creates an associative array that stores student information:

\begin{minted}[]{php}
<?=
 $student = [
     "name" => "Alice",
     "age" => 20,
     "course" => "Computer Science"
 ];
\end{minted}

In this example, the \texttt{student} array has three elements, each
represented by a key-value pair. The key is the label for the data, and
the value is the actual data itself.

\textbf{Accessing Elements in Associative Arrays}

To access an element in an associative array, you use the key instead of
the index. For example, the following code accesses the student's name
from the \texttt{student} array:

\begin{minted}[]{php}
<?=
 $firstName = $student["name"];
 echo "Student's name: $firstName";
\end{minted}

\textbf{Adding Elements to Associative Arrays}

To add new elements to an associative array, you simply add new
key-value pairs. For instance, the following code adds the student's GPA
to the \texttt{student} array:

\begin{minted}[]{php}
<?=
 $student["gpa"] = 3.9;
 echo "Student's GPA: " . $student["gpa"];
\end{minted}

\textbf{Removing Elements from Associative Arrays}

To remove an element from an associative array, you use the \texttt{unset()}
function and provide the key of the element to be removed. For example,
the following code removes the student's age from the \texttt{student} array:

\begin{minted}[]{php}
<?=
 unset($student["age"]);
 echo "Student's age: " . $student["age"];
\end{minted}

Associative arrays are a versatile and efficient data structure for
storing and retrieving information based on meaningful labels. They are
commonly used in various applications, such as user profiles, product
catalogs, and configuration settings.

\subsection{Here's a list of functions for working with arrays in PHP grouped by}
\label{sec:org0ea5afb}
functionality:

\subsubsection{Array creation and initialization:*}
\label{sec:orgd30115d}
\begin{itemize}
\item \texttt{array()}: Creates an empty array or initializes an array with a set
of values.
\item \texttt{count()}: Returns thenumber of elements in an array.
\item \texttt{in\_array()}: Checks whether a specific value exists in an array.
\end{itemize}

\subsubsection{Array manipulation and modification:*}
\label{sec:org6c55f78}
\begin{itemize}
\item \texttt{array\_push()}: Adds one or more elements to the end of an array.
\item \texttt{array\_pop()}: Removes and returns the last element of an array.
\item \texttt{array\_shift()}: Removes and returns the first element of an array.
\item \texttt{array\_unshift()}: Adds one or more elements to the beginning of an array.
\item \texttt{array\_splice()}: Removes and replaces elements from an array.
\item \texttt{isset()}: Checks whether a specific key exists in an array.
\item \texttt{empty()}: Checks whether an array isempty.
\item \texttt{unset()}: Removes one or more elements from an array.
\end{itemize}

\subsubsection{Array iteration and searching:*}
\label{sec:org04383c9}
\begin{itemize}
\item \texttt{foreach()}: Iterates over the elements of an array.
\item \texttt{array\_keys()}: Returns an array of all the keys in an array.
\item \texttt{array\_values()}: Returns an array of all the values in an array.
\item \texttt{array\_search()}: Finds the first occurrence of a specific value in an array.
\end{itemize}

\subsubsection{Array sorting and filtering:*}
\label{sec:orge676bb2}
\begin{itemize}
\item \texttt{sort()}: Sorts an array in ascending order.
\item \texttt{rsort()}: Sorts an array in descending order.
\item \texttt{asort()}: Sorts an associative array by value.
\item \texttt{arsort()}: Sorts an associative array by value in descending order.
\item \texttt{ksort()}: Sorts an associative array by key.
\item \texttt{krsort()}: Sorts an associative array by key in descending order.
\item \texttt{array\_filter()}: Filters an array based on a callback function.
\end{itemize}

\subsubsection{Array merging and combining:*}
\label{sec:orgc3fad54}
\begin{itemize}
\item \texttt{array\_merge()}: Merges two or more arrays into one.
\item \texttt{array\_slice()}: Extracts a portion of an array to create a new array.
\item \texttt{array\_unique()}: Removes duplicate values from an array.
\item \texttt{array\_fill()}: Creates a new array with a specified value and a range of keys.
\end{itemize}

\subsubsection{Array accessing and manipulating elements:*}
\label{sec:org330c103}
\begin{itemize}
\item \texttt{array\_keys()}: Returns an array of all the keys in an array.
\item \texttt{array\_values()}: Returns an array of all the values in an array.
\item \texttt{array\_chunk()}: Divides an array into smaller subarrays.
\item \texttt{array\_combine()}: Creates an associative array from two arrays of keys and values.
\item \texttt{array\_map()}: Applies a callback function to each element in an array.
\item \texttt{array\_reduce()}: Reduces an array to a single value using a callback function.
\end{itemize}

These functions provide a comprehensive set of tools for working with
arrays in PHP, covering a wide range of tasks from creating and
manipulating arrays to searching, sorting, and filtering them.
\end{document}
